\documentclass[12pt]{article}
\usepackage{amsmath,amssymb}

\begin{document}
\section*{{2020 AIME I Problems/Problem 14}}
Source: AoPS Wiki

\subsection*{Solution}
Either $P(3) = P(4)$ or not. We first see that if $P(3) = P(4)$ it's easy to obtain by Vieta's that $(a+b)^2 = 49$ . Now, take $P(3) \neq P(4)$ and WLOG $P(3) = P(a), P(4) = P(b)$ . Now, consider the parabola formed by the graph of $P$ . It has vertex $\frac{3+a}{2}$ . Now, say that $P(x) = x^2 - (3+a)x + c$ . We note $P(3)P(4) = c = P(3)\left(4 - 4a + \frac{8a - 1}{2}\right) \implies a = \frac{7P(3) + 1}{8}$ . Now, we note $P(4) = \frac{7}{2}$ by plugging in again. Now, it's easy to find that $a = -2.5, b = -3.5$ , yielding a value of $36$ . Finally, we add $49 + 36 = \boxed{085}$ . ~awang11, charmander3333 Remark : We know that $c=\frac{8a-1}{2}$ from $P(3)+P(4)=3+a$ .

\end{document}
